% ============================================================
%  Estils TikZ comuns per als diagrames del llibre.
%
%  Els diagrames de teoria de categories repeteixen molt codi de
%  puntes de fletxa i de matrius de nodes. Aquests estils permeten
%  escriure els diagrames nous de manera compacta i canviar
%  l'estètica de tots des d'un sol lloc:
%
%     \begin{tikzpicture}
%       \matrix (m) [cat matrix] { A & B \\ C & D \\ };
%       \draw[cat] (m-1-1) -- node[above] {$f$} (m-1-2);
%     \end{tikzpicture}
%
%  Els diagrames antics, escrits amb -{Stealth[scale=1.1]} i amb
%  els paràmetres de matriu explícits, continuen funcionant sense
%  canvis; aquests estils són per a diagrames nous i per a una
%  homogeneïtzació progressiva.
% ============================================================

\tikzset{
  % Fletxa estàndard d'un morfisme.
  cat/.style={-{Stealth[scale=1.1]}},
  % Fletxa de monomorfisme (cua d'ham).
  cat mono/.style={{Hooks[right]}-{Stealth[scale=1.1]}},
  % Fletxa d'epimorfisme (doble punta).
  cat epi/.style={-{Stealth[scale=1.1]Stealth[scale=1.1]}},
  % Fletxa de transformació natural (doble).
  cat nat/.style={-{Implies},double,double distance=1.5pt},
  % Matriu de nodes amb la separació habitual del llibre.
  cat matrix/.style={matrix of math nodes, row sep=2.6em, column sep=3.2em},
}